%! TEX root = ../ag-19-20.tex

\chapter{Examen 2018--2019}

\textbf{Numărul 1}

1. În spațiul vectorial $ \RR^3 $ se consideră mulțimea:
\[
  V = \{ (a, b, c) \in \RR^3 \mid a - 2b + 3c = 0 \}.
\]

Să se determine $V^\perp$ relativ la produsul scalar euclidian și o
bază ortonormată în $ \RR^3 $ diferită de baza canonică.

\vspace{1cm}

2. Să se aducă la forma canonică și să se reprezinte grafic conica:
\[
  f(x, y) = 9x^2 + 4xy - 3y^2 + 4x - 14y - 7 = 0.
\]


\vspace{1cm}

3. Să se rezolve ecuațiile diferențiale:
\begin{enumerate}[(a)]
  \item $ xy' - y + 2x^2y^2 = 0 $;
  \item $ y = x(1 + y') + (y')^2 $;
  \item $ (3x^2 + 2) y' + xy^2 = 0 $;
  \item $ (2x - y + 2) dx + (-x + y + 1) dy = 0 $.
\end{enumerate}

\vspace{1cm}

4. Să se rezolve sistemul:
\[
  \begin{cases}
    x' &= -3x - y + e^t \\
    y' &= x - y
  \end{cases}
\]
cu condițiile inițiale $ x(0) = 1, y(0) = -1 $ și să se calculeze
$ e^{At}$, unde $ A $ este matricea sistemului.

\vspace{1cm}

5.(a) Să se determine ecuația conului cu vîrful în $ V(2, -1, 0) $ și
determinat de curba\footnotemark:
\[
  (\gamma) : \begin{cases}
    y^2 + z^2 &= 2 \\
    x &= 1
  \end{cases}
\]

\footnotetext{Pentru ecuația conului și a altor cuadrice nestudiate
    la seminar, consultați notițele Prof.\ Niță sau, de exemplu,
\href{https://www.deliu.ro/pluginfile.php/109/mod_resource/content/1/cap02.pdf}{acest document}.}

(b) Se consideră spațiul euclidian $ \RR_4[X] $ și un endomorfism
$ f : \RR_4[X] \to \RR_4[X] $. Să se arate că:
\[
  \dr{Im}f \simeq (\dr{Ker}f^\ast)^\perp.
\]



\textbf{Numărul 2}

1. În spațiul euclidian $ \RR^3 $ se consideră mulțimea:
\[
  V = \{(a, b, c) \in \RR^3 \mid 2a - b + 3c = 0 \}.
\]

Să se determine $ V^\perp $ și o bază ortonormată în $ \RR^3 $, diferită
de baza canonică.

\vspace{1cm}


2. Să se aducă la forma canonică și să se reprezinte grafic conica:
\[
  f(x, y) = 9x^2 + 4xy + 3y^2 + 16x + 12y - 36 = 0.
\]

\vspace{1cm}

3. Să se rezvole ecuațiile diferențiale:
\begin{enumerate}[(a)]
  \item $ xy' = 4y + x^2 \sqrt{y} $;
  \item $ yy' = 2x(y')^2 + 1 $;
  \item $ (x^2 + 1)y' + 2xy^2 = 0 $;
  \item $ (2x + y + 1)dx + (x + y - 1) dy = 0 $.
\end{enumerate}


\vspace{1cm}

4. Să se rezolve sistemul:
\[
  \begin{cases}
    x' &= 2x + y \\
    y' &= -x + 4y + 27t
  \end{cases}
\]
cu valorile inițiale $ x(0) = -1, y(0) = 1 $ și să se calculeze
$ e^{At} $, unde $ A $ este matricea sistemului.

\vspace{1cm}

5.(a) Să se determine ecuația cilindrului cu generatoarele paralele
cu vectorul $ \vec{v} = \vec{i} - \vec{j} $ și curba directoare:
\[
  (\gamma): %
  \begin{cases}
    x^2 - y^2 + z^2 &= 0 \\
    x &= 1
  \end{cases}
\]

(b) Se consideră spațiul euclidian $ \RR_4[X] $ și $f : \RR_4[X] \to \RR_4[X] $
un endomorfism. Să se arate că:
\[
  \dr{Ker}f \simeq (\dr{Im}f^\ast)^\perp.
\]


%%% Local Variables:
%%% mode: latex
%%% TeX-master: "../ag-20-21"
%%% End:
